% paper/sections/appendix_c_ccd.tex
%
% 付録C: CCD 法の係数導出と実装（旧付録B+C を統合）
%   — §4/§8 対応

\section{CCD 法の係数導出と実装}
\label{app:ccd_supplement}

%% C.1 CCD 係数の Taylor 展開導出（内点＋境界）
\subsection{CCD 係数の Taylor 展開導出}
\label{app:ccd_coef_all}
\subsubsection{内点係数の Taylor 展開による導出}
\label{app:ccd_coef_derivation}
% ═══════════════════════════════════════════════════════════════════════════════

本付録は §\ref{sec:ccd_te_I} で述べた CCD スキームの係数
$(\alpha_1, a_1, b_1)$ および $(\beta_2, a_2, b_2)$ を，
Taylor 展開による打ち切り誤差解析から厳密に導出する．
各 Equation の RHS を格子点 $i$ 周りで展開し，
LHS との係数比較により3元連立方程式を立て，代数的に解く．

\subsubsection{Equation-I の係数導出}
\label{app:ccd_coef_derivation_I}

Equation-I（式~\eqref{eq:CCD_I}）の各項を $f$ の Taylor 展開で表す．
添字 $i\pm1$ の一次微分値と二次微分値は：
\begin{align}
  f'_{i\pm1} &= f'_i \pm hf''_i + \frac{h^2}{2}f^{(3)}_i \pm \frac{h^3}{6}f^{(4)}_i
               + \frac{h^4}{24}f^{(5)}_i \pm \frac{h^5}{120}f^{(6)}_i
               + \frac{h^6}{720}f^{(7)}_i + \Ord{h^7} \\
  f''_{i\pm1} &= f''_i \pm hf^{(3)}_i + \frac{h^2}{2}f^{(4)}_i \pm \frac{h^3}{6}f^{(5)}_i
                + \frac{h^4}{24}f^{(6)}_i \pm \frac{h^5}{120}f^{(7)}_i + \Ord{h^6}
\end{align}
関数値については：
\begin{align}
  f_{i+1} - f_{i-1} &= 2hf'_i + \frac{h^3}{3}f^{(3)}_i + \frac{h^5}{60}f^{(5)}_i
                       + \frac{h^7}{2520}f^{(7)}_i + \Ord{h^8}
\end{align}

\subsubsection*{左辺（LHS）の展開}

\begin{align}
  \text{LHS} &= \alpha_1 f'_{i-1} + f'_i + \alpha_1 f'_{i+1}
  = (1 + 2\alpha_1)f'_i + \alpha_1 h^2 f^{(3)}_i
    + \frac{\alpha_1}{12}h^4 f^{(5)}_i + \frac{\alpha_1}{360}h^6 f^{(7)}_i + \Ord{h^7}
\end{align}

\subsubsection*{右辺（RHS）の展開}

\begin{align}
  \frac{a_1}{h}(f_{i+1}-f_{i-1})
  &= 2a_1 f'_i + \frac{a_1}{3}h^2 f^{(3)}_i + \frac{a_1}{60}h^4 f^{(5)}_i
     + \frac{a_1}{2520}h^6 f^{(7)}_i + \Ord{h^7} \\
  b_1 h\bigl(f''_{i+1}-f''_{i-1}\bigr)
  &= 2b_1 h^2 f^{(3)}_i + \frac{b_1}{3}h^4 f^{(5)}_i
     + \frac{b_1}{60}h^6 f^{(7)}_i + \Ord{h^7}
\end{align}

合計：
\begin{equation}
  \text{RHS} = 2a_1 f'_i
  + \Bigl(\frac{a_1}{3} + 2b_1\Bigr)h^2 f^{(3)}_i
  + \Bigl(\frac{a_1}{60} + \frac{b_1}{3}\Bigr)h^4 f^{(5)}_i
  + \Bigl(\frac{a_1}{2520} + \frac{b_1}{60}\Bigr)h^6 f^{(7)}_i + \Ord{h^7}
\end{equation}

\subsubsection*{係数比較と連立方程式}

$\Ord{h^6}$ まで LHS $=$ RHS が成立するには，
$f'_i$，$h^2 f^{(3)}_i$，$h^4 f^{(5)}_i$ の各係数が一致しなければならない：

\begin{align}
  f'_i       &:\quad 1 + 2\alpha_1 = 2a_1 \label{eq:app_I_eq1} \\
  h^2 f^{(3)}_i &:\quad \alpha_1 = \frac{a_1}{3} + 2b_1 \label{eq:app_I_eq2} \\
  h^4 f^{(5)}_i &:\quad \frac{\alpha_1}{12} = \frac{a_1}{60} + \frac{b_1}{3} \label{eq:app_I_eq3}
\end{align}

\subsubsection*{解法}

式~\eqref{eq:app_I_eq1} より $a_1 = (1+2\alpha_1)/2$．
式~\eqref{eq:app_I_eq2} より $b_1 = (\alpha_1 - a_1/3)/2$：
\[
  b_1 = \frac{1}{2}\!\left(\alpha_1 - \frac{1+2\alpha_1}{6}\right)
      = \frac{1}{2}\cdot\frac{6\alpha_1 - 1 - 2\alpha_1}{6}
      = \frac{4\alpha_1 - 1}{12}
\]
$a_1$ と $b_1$ を式~\eqref{eq:app_I_eq3} に代入：
\[
  \frac{\alpha_1}{12}
  = \frac{1+2\alpha_1}{120} + \frac{4\alpha_1-1}{36}
\]
両辺に $360$ を掛けて整理：
\[
  30\alpha_1 = 3(1+2\alpha_1) + 10(4\alpha_1-1)
             = 3 + 6\alpha_1 + 40\alpha_1 - 10
             = 46\alpha_1 - 7
\]
\[
  -16\alpha_1 = -7
  \quad\Longrightarrow\quad
  \alpha_1 = \frac{7}{16}
\]

以上より：
\begin{alignat}{3}
  \alpha_1 &= \frac{7}{16}, \qquad
  &a_1 &= \frac{1+2\cdot(7/16)}{2} = \frac{1 + 7/8}{2} = \frac{15}{16}, \qquad
  &b_1 &= \frac{4\cdot(7/16)-1}{12} = \frac{7/4-1}{12} = \frac{3/4}{12} = \frac{1}{16}
\end{alignat}

\subsubsection*{打ち切り誤差}

$h^6 f^{(7)}_i$ の LHS 係数 $-$ RHS 係数：
\[
  TE_{\mathrm{I}} = \left(\frac{\alpha_1}{360} - \frac{a_1}{2520} - \frac{b_1}{60}\right)h^6 f^{(7)}_i
\]
数値代入：
\begin{align*}
  &\frac{7/16}{360} - \frac{15/16}{2520} - \frac{1/16}{60}
  = \frac{1}{16}\!\left(\frac{7}{360} - \frac{15}{2520} - \frac{1}{60}\right) \\
  &= \frac{1}{16}\!\left(\frac{49}{2520} - \frac{15}{2520} - \frac{42}{2520}\right)
  = \frac{1}{16}\cdot\frac{-8}{2520}
  = -\frac{1}{5040} = -\frac{1}{7!}
\end{align*}
したがって
\[
  TE_{\mathrm{I}} = -\frac{1}{7!}\,h^6 f^{(7)}_i = \Ord{h^6}
\]
が成立し，スキームが $\Ord{h^6}$ 精度であることが確認された．\qed

\subsubsection{Equation-II の係数導出}
\label{app:ccd_coef_derivation_II}

Equation-II（式~\eqref{eq:CCD_II}）の各項を展開する．
二次微分値と関数値の差分については：
\begin{align}
  f''_{i\pm1} &= f''_i \pm hf^{(3)}_i + \frac{h^2}{2}f^{(4)}_i \pm \frac{h^3}{6}f^{(5)}_i
                + \frac{h^4}{24}f^{(6)}_i \pm \frac{h^5}{120}f^{(7)}_i
                + \frac{h^6}{720}f^{(8)}_i + \Ord{h^7} \\
  f_{i-1} - 2f_i + f_{i+1} &= h^2 f''_i + \frac{h^4}{12}f^{(4)}_i
                               + \frac{h^6}{360}f^{(6)}_i
                               + \frac{h^8}{20160}f^{(8)}_i + \Ord{h^8} \\
  f'_{i+1} - f'_{i-1} &= 2hf''_i + \frac{h^3}{3}f^{(4)}_i
                         + \frac{h^5}{60}f^{(6)}_i
                         + \frac{h^7}{2520}f^{(8)}_i + \Ord{h^8}
\end{align}

\subsubsection*{左辺（LHS）の展開}

\begin{equation}
  \text{LHS} = \beta_2 f''_{i-1} + f''_i + \beta_2 f''_{i+1}
  = (1+2\beta_2)f''_i + \beta_2 h^2 f^{(4)}_i
    + \frac{\beta_2}{12}h^4 f^{(6)}_i + \frac{\beta_2}{360}h^6 f^{(8)}_i + \Ord{h^7}
\end{equation}

\subsubsection*{右辺（RHS）の展開}

\begin{align}
  \frac{a_2}{h^2}(f_{i-1}-2f_i+f_{i+1})
  &= a_2 f''_i + \frac{a_2}{12}h^2 f^{(4)}_i + \frac{a_2}{360}h^4 f^{(6)}_i
     + \frac{a_2}{20160}h^6 f^{(8)}_i + \Ord{h^7} \\
  \frac{b_2}{h}(f'_{i+1}-f'_{i-1})
  &= 2b_2 f''_i + \frac{b_2}{3}h^2 f^{(4)}_i + \frac{b_2}{60}h^4 f^{(6)}_i
     + \frac{b_2}{2520}h^6 f^{(8)}_i + \Ord{h^7}
\end{align}

合計：
\begin{equation}
  \text{RHS} = (a_2+2b_2)f''_i
  + \Bigl(\frac{a_2}{12}+\frac{b_2}{3}\Bigr)h^2 f^{(4)}_i
  + \Bigl(\frac{a_2}{360}+\frac{b_2}{60}\Bigr)h^4 f^{(6)}_i
  + \Bigl(\frac{a_2}{20160}+\frac{b_2}{2520}\Bigr)h^6 f^{(8)}_i + \Ord{h^7}
\end{equation}

\subsubsection*{係数比較と連立方程式}

$f''_i$，$h^2 f^{(4)}_i$，$h^4 f^{(6)}_i$ の係数を比較する：
\begin{align}
  f''_i        &:\quad (1+2\beta_2) - (a_2+2b_2) = 0 \label{eq:app_II_eq1} \\
  h^2 f^{(4)}_i &:\quad \beta_2 - \Bigl(\frac{a_2}{12}+\frac{b_2}{3}\Bigr) = 0 \label{eq:app_II_eq2} \\
  h^4 f^{(6)}_i &:\quad \frac{\beta_2}{12} - \Bigl(\frac{a_2}{360}+\frac{b_2}{60}\Bigr) = 0 \label{eq:app_II_eq3}
\end{align}

\subsubsection*{解法}

式~\eqref{eq:app_II_eq2} および \eqref{eq:app_II_eq3} から $\beta_2$ を消去する．
式~\eqref{eq:app_II_eq2} の12倍 $=$ 式~\eqref{eq:app_II_eq3} の左辺に代入：
\[
  \frac{a_2}{12} + \frac{b_2}{3} = \frac{a_2}{30} + \frac{b_2}{5}
\]
整理：
\[
  a_2\!\left(\frac{1}{12}-\frac{1}{30}\right) = b_2\!\left(\frac{1}{5}-\frac{1}{3}\right)
  \quad\Longrightarrow\quad
  a_2\cdot\frac{3}{60} = b_2\cdot\!\left(-\frac{2}{15}\right)
  \quad\Longrightarrow\quad
  b_2 = -\frac{3a_2}{8}
\]
式~\eqref{eq:app_II_eq1} に代入：
\[
  1 + 2\beta_2 = a_2 + 2b_2 = a_2 - \frac{3a_2}{4} = \frac{a_2}{4}
  \quad\Longrightarrow\quad
  \beta_2 = \frac{a_2}{8} - \frac{1}{2}
\]
これを式~\eqref{eq:app_II_eq2} に代入：
\[
  \frac{a_2}{8} - \frac{1}{2} = \frac{a_2}{12} + \frac{1}{3}\cdot\!\left(-\frac{3a_2}{8}\right)
  = \frac{a_2}{12} - \frac{a_2}{8}
  = -\frac{a_2}{24}
\]
\[
  \frac{a_2}{8} + \frac{a_2}{24} = \frac{1}{2}
  \quad\Longrightarrow\quad
  \frac{4a_2}{24} = \frac{1}{2}
  \quad\Longrightarrow\quad
  a_2 = 3
\]

以上より：
\begin{alignat}{3}
  a_2 &= 3, \qquad
  &b_2 &= -\frac{3\times3}{8} = -\frac{9}{8}, \qquad
  &\beta_2 &= \frac{3}{8} - \frac{1}{2} = -\frac{1}{8}
\end{alignat}

\subsubsection*{打ち切り誤差}

$h^6 f^{(8)}_i$ の LHS 係数 $-$ RHS 係数：
\[
  TE_{\mathrm{II}} = \left(\frac{\beta_2}{360} - \frac{a_2}{20160} - \frac{b_2}{2520}\right)h^6 f^{(8)}_i
\]
数値代入：
\[
  \frac{-1/8}{360} - \frac{3}{20160} - \frac{-9/8}{2520}
  = -\frac{1}{2880} - \frac{3}{20160} + \frac{9}{20160}
  = -\frac{7}{20160} + \frac{6}{20160}
  = -\frac{1}{20160}
\]
したがって
\[
  TE_{\mathrm{II}} = -\frac{1}{20160}\,h^6 f^{(8)}_i = \Ord{h^6}
\]
が成立し，Equation-II も $\Ord{h^6}$ 精度であることが確認された．\qed

\medskip
\noindent
主要打ち切り誤差係数 $-1/5040 = -1/7!$（Eq-I）と $-1/20160$（Eq-II）はいずれも
7点陽的スキームの打ち切り誤差係数と同オーダーであり，
3点ステンシルで7点相当の精度を実現していることを示している．

% ═══════════════════════════════════════════════════════════════════════════════
 %% C.1.1 内点係数
\input{sections/appendix_c1_2_ccd_boundary_coefficients} %% C.1.2 境界スキーム係数

%% C.2 境界処理の代替実装（ゴーストセル法＋周期BC）
\subsection{境界処理の代替実装}
\label{app:ccd_boundary_alt}
\input{sections/appendix_c2_1_ccd_ghost_cells} %% C.2.1 ゴーストセル法
\subsubsection{周期境界条件：ブロック循環 CCD 行列法}
\label{app:ccd_periodic}
% ═══════════════════════════════════════════════════════════════════════════════

\begin{tcolorbox}[warnbox, title={片側スキームと周期 BC の相性問題}]
§\ref{sec:ccd_bc} の片側差分境界スキームは，\textbf{非周期的な両端固定境界}（壁面・Dirichlet・Neumann）
に対して設計されている．周期境界条件（$f_0 = f_N$）に対してこのスキームをそのまま適用すると，
境界節点 $i=0,\, N$ の微分値に浮動小数点丸め起因の $\mathcal{O}(\epsilon_{\mathrm{mach}}/h)$ 程度の
誤差が残留する．

\medskip
\textbf{不安定化の機構．}
境界節点の誤差が AB2 時間積分スキームによって増幅される現象は，
Gustafsson \cite{Gustafsson1981} の境界安定性問題として知られている．
AB2 法の安定領域は純虚数軸上で $|\mathrm{Im}(\hat{\lambda})| \lesssim 0.47$ に制限されるが，
片側スキームが導入する非対称な境界行は実部正の偽固有値を生成し，AB2 の安定領域を逸脱する．
実際，表~\ref{tab:uniform_flow_comparison} に示すように，
一様流試験（$u=1$，周期 BC）でこのスキームを誤用すると
1 ステップあたり約 16 倍の誤差増幅が生じ，$t \approx 0.07$ でブローアップする．

\medskip
\textbf{根本原因：}
節点 $i=0$ の内点ステンシルが必要とする左隣り $f_{-1}$ は，
周期境界では $f_{N-1}$（折り返し）であるが，
片側スキームは代わりに $\{f_0, f_1, f_2, f_3\}$ を用いる—これは物理的に誤った構成である．
\end{tcolorbox}

\subsubsection{正しい定式化：全節点に内点スキームを適用}

周期境界条件下では，$N$ 個の独立節点（$i = 0, \ldots, N-1$；節点 $N$ は節点 0 の周期像）が
\textbf{すべて内点スキーム（式~\eqref{eq:ccd_block}）を折り返し添字で使用可能}である：
%
\begin{equation}
  \mathbf{A}_L\,\bm{v}_{(i-1)\bmod N}
  + \mathbf{B}\,\bm{v}_i
  + \mathbf{A}_R\,\bm{v}_{(i+1)\bmod N}
  = \mathbf{r}_i, \qquad i = 0, 1, \ldots, N-1
  \label{eq:ccd_periodic}
\end{equation}
%
右辺は折り返し添字で評価する：
%
\begin{equation}
  \mathbf{r}_i =
  \begin{bmatrix}
    \dfrac{a_1}{h}\bigl(f_{(i+1)\bmod N} - f_{(i-1)\bmod N}\bigr) \\[8pt]
    \dfrac{a_2}{h^2}\bigl(f_{(i-1)\bmod N} - 2f_i + f_{(i+1)\bmod N}\bigr)
  \end{bmatrix}
  \label{eq:rhs_periodic}
\end{equation}
%
求解後，周期像節点の微分値は $\bm{v}_N = \bm{v}_0$ とおく（同一物理点）．

\subsubsection{全体行列：\texorpdfstring{$2N\times 2N$}{2N x 2N} ブロック循環行列}

$N$ 式（各 $2\times 2$ ブロック）をまとめた全体行列は
\textbf{ブロック循環（block-circulant）行列}となる：
%
\begin{equation}
  \mathbf{K}_{\mathrm{circ}} =
  \begin{bmatrix}
    \mathbf{B}   & \mathbf{A}_R & \mathbf{0}   & \cdots & \mathbf{0}   & \mathbf{A}_L \\
    \mathbf{A}_L & \mathbf{B}   & \mathbf{A}_R & \cdots & \mathbf{0}   & \mathbf{0}   \\
    \mathbf{0}   & \mathbf{A}_L & \mathbf{B}   & \ddots & \vdots       & \vdots       \\
    \vdots       &              & \ddots        & \ddots & \mathbf{A}_R & \mathbf{0}   \\
    \mathbf{0}   & \cdots       & \mathbf{0}   & \mathbf{A}_L & \mathbf{B} & \mathbf{A}_R \\
    \mathbf{A}_R & \mathbf{0}   & \cdots & \mathbf{0}   & \mathbf{A}_L & \mathbf{B}
  \end{bmatrix}
  \in \mathbb{R}^{2N\times 2N}
  \label{eq:ccd_circulant}
\end{equation}
%
ここで $\mathbf{A}_L,\, \mathbf{A}_R,\, \mathbf{B}$ は式~\eqref{eq:ccd_block} の内点ブロックと完全に同一であり，
\textbf{境界専用ブロックは一切不要}となる．
ブロック三重対角行列（式~\eqref{eq:ccd_global}）との違いは隅の $\mathbf{A}_L$（右上：$i=0$ 行の折り返し）と $\mathbf{A}_R$（左下：$i=N{-}1$ 行の折り返し）の2ブロックのみである．

\begin{tcolorbox}[defbox, title={ブロック循環行列の正則性}]
$\mathbf{K}_{\mathrm{circ}}$ が正則であることを示す．

定数モード $\bm{v}_i = [c,\, 0]^\top$ を代入すると各行は
$(\mathbf{A}_L + \mathbf{B} + \mathbf{A}_R)[c,\, 0]^\top = [(1+2\alpha_1)c,\, 0]^\top$ となり，
$1+2\alpha_1 = 1 + 7/8 = 15/8 \neq 0$ だから零固有値にならない．

定数モード $\bm{v}_i = [0,\, c]^\top$ は $(1+2\beta_2) = 3/4 \neq 0$ から同様に非退化．

一般の波数 $k=1,\ldots,N-1$ の Fourier モード $\bm{v}_i = e^{2\pi\mathrm{i}ki/N}\bm{w}$ に対する固有値は，
$\hat{\mathbf{K}}_k = e^{-2\pi\mathrm{i}k/N}\mathbf{A}_L + \mathbf{B} + e^{+2\pi\mathrm{i}k/N}\mathbf{A}_R$ の固有値に対応し，
CCD スキームの修正波数解析（§\ref{sec:ccd_def}）から全波数で非零である．
ゆえに $\mathbf{K}_{\mathrm{circ}}$ は正則である．
\end{tcolorbox}

\subsubsection{数値解法：密行列 LU 分解}

$N \leq 512$ 程度では，$2N\times 2N$ 行列の密行列 LU 分解（前計算コスト $\mathcal{O}(N^3)$，
$N=32$ で $64\times64$，$N=128$ で $256\times256$）を初期化時に1回だけ実行し，
タイムステップ毎に \texttt{lu\_solve} によるバッチ後退代入（コスト $\mathcal{O}(N^2 B)$，$B$ はバッチサイズ）を適用する．

\begin{tcolorbox}[algbox, title={周期 CCD 解法の手順（実装概要）}]
\textbf{初期化（$\mathcal{O}(N^3)$，1回）：}
\begin{enumerate}
  \item $2N\times 2N$ ブロック循環行列 $\mathbf{K}_{\mathrm{circ}}$（式~\eqref{eq:ccd_circulant}）を組み立てる．
  \item $(\mathbf{L},\,\mathbf{U},\,\mathbf{P}) = \mathrm{lu}(\mathbf{K}_{\mathrm{circ}})$ を計算・保存する．
\end{enumerate}
\textbf{求解（$\mathcal{O}(N^2 B)$，タイムステップ毎）：}
\begin{enumerate}
  \item 折り返し添字で右辺 $\mathbf{R} \in \mathbb{R}^{2N\times B}$ を構築する（式~\eqref{eq:rhs_periodic}）．
  \item $\mathbf{K}_{\mathrm{circ}}\,\mathbf{V} = \mathbf{R}$ を LU 後退代入で解く．
  \item $[f'_i,\, f''_i]^\top = \mathbf{V}_{2i:2i+2}$ を取り出し，$\bm{v}_N = \bm{v}_0$ とする．
\end{enumerate}
\end{tcolorbox}

\subsubsection{検証：一様流試験}
\label{sec:verify_uniform_flow}

実装の正確性を確認するために，\textbf{一様流試験}（$u\equiv1,\; v\equiv0,\; p\equiv0$；
$32\times32$ 格子，$Re=10^6$，$We=Fr=10^{10}$，$\rho_r=\mu_r=1$，CFL$=0.3$，周期 BC）を実施した．
理論上，全ての項（対流・粘性・IPC 投影・面フラックス補正）が恒等的にゼロとなるため，
速度・圧力は全時刻で初期値のまま保持されなければならない．

\begin{table}[htbp]
\centering
\caption{一様流試験の結果比較：周期 CCD と片側スキーム誤適用}
\label{tab:uniform_flow_comparison}
\begin{tabular}{lccccc}
\toprule
スキーム & ステップ & $t$ & $\|u-1\|_\infty$ & $\|v\|_\infty$ & $\|p\|_\infty$ \\
\midrule
\multirow{3}{*}{周期 CCD（正）}
  &  1 & 0.009 & $0$              & $9\times10^{-23}$ & $0$ \\
  & 10 & 0.094 & $0$              & $9\times10^{-22}$ & $0$ \\
  & 20 & 0.188 & $0$              & $2\times10^{-21}$ & $0$ \\
\midrule
\multirow{3}{*}{片側スキーム（誤）}
  &  1 & 0.009 & $2.5\times10^{-9}$  & $2.5\times10^{-9}$  & $3.3\times10^{-8}$ \\
  &  7 & 0.066 & $2.5\times10^{-3}$  & $2.6\times10^{-3}$  & $2.8\times10^{-1}$ \\
  & 10 & 0.094 & $4.7\times10^{0}$   & $5.0\times10^{0}$   & $1.4\times10^{3}$  \\
\bottomrule
\end{tabular}
\end{table}

\noindent
周期 CCD では $\|u-1\|_\infty$ および $\|p\|_\infty$ は全ステップで浮動小数点表現の最小単位（マシンゼロ）であり，
$\|v\|_\infty$ は 20 ステップで $\mathcal{O}(10^{-21})$ と線形増加する（時間積分の丸め累積のみ）．
一方，片側スキームを誤用した場合は 1 ステップあたり約 16 倍の誤差増幅が生じ，$t\approx0.07$ でブローアップした．
 %% C.2.2 周期境界条件

%% C.3 精度解析と楕円型ソルバー（混合微分精度＋PPEソルバー構造＋Kronecker積）
\subsection{精度解析と楕円型ソルバー}
\label{app:ccd_analysis}
\input{sections/appendix_c3_1_ccd_mixed_derivative} %% C.3.1 混合微分精度解析
\input{sections/appendix_c3_2_ccd_elliptic_solver}  %% C.3.2 楕円型ソルバー・Kronecker積・Neumann BC
